\section{Critérios e Métodos de Priorização}
\label{sec:criterios}

\subsection{Critérios de priorização}
Apresentar os critérios mais usados na literatura, idealmente numa pequena tabela com as referências que usam cada um.
Os critérios mais recorrentes incluem a criticidade do equipamento, o impacto na segurança, o impacto na produção, o custo da intervenção, a probabilidade de falha, o histórico de falhas, a redundância e a disponibilidade de recursos. Rastogi et al.~\cite{rastogi2024} identificam nove fatores de avaliação, embora apenas quatro sejam efetivamente usados no seu modelo.

\subsection{Análise de criticidade e FMEA}
Matrizes de criticidade, FMEA e o Risk Priority Number (RPN = severidade $\times$ ocorrência $\times$ deteção); limitações conhecidas do RPN.

\subsection{Métodos de decisão multicritério}
Para cada método, uma ou duas frases sobre como funciona e um exemplo de aplicação à manutenção:
\begin{itemize}
  \item \textbf{Definição de pesos:} AHP, BWM.
  \item \textbf{Ordenação de alternativas:} soma ponderada, TOPSIS, VIKOR, PROMETHEE, ELECTRE.
  \item \textbf{Variantes fuzzy:} tratamento da incerteza nos julgamentos dos especialistas.
\end{itemize}

\subsection{Otimização e escalonamento}
Programação matemática, meta-heurísticas (ex.: PSO, algoritmos genéticos) e escalonamento com restrições de recursos e janelas de paragem.